%!TEX root = ../main.tex
% =============================================================================
%  七、问题 4：波动电价下的风险决策模型      建议 3–4 页
%  评审在问：波动电价怎么处理的？
%  本节必须出现：① 电价过程的四项检验  ② 两层建模（分时骨架 + 随机残差）
%                ③ 联合情景（二维分层取代表历史日），必须保留 rho(p, L)
%                ④ 波动电价下的 LP 形式（双重不确定性）
%                ⑤ 风险—费用前沿 + 真实轨迹回测（口径 C-11）
%  数据来源：notes/数值台账.md、intermediate/paper_numbers.json（键 电价检验/Q4_2/Q4_3）、
%            intermediate/q4_price.json、intermediate/q4_cvar.json
%  口径编号：按 notes/写作交底.md §6.1 的统一分配，CVaR 用 C-11
%  图文件名：fig08_q4_price.pdf（电价骨架与残差自相关）、fig09_cvar_frontier.pdf（前沿与回测）
% =============================================================================

\section{问题 4：波动电价下的风险决策模型}

\subsection{电价过程的统计检验}

问题 4 中电价与光伏出力同时随机；若把附件 4 的电价序列直接当作"另一组已知价格"代入模型，等于假设
\(0{:}00\) 就知道全天的实际电价，与题意矛盾。故本文先把电价当作\textbf{随机过程}检验并建模
（表~\ref{tab:price-test}）：检验在 2025 年全年（\(52\,560\) 个观测）上进行，费用结果仍只在 334 天上报告。



四项检验的建模含义见表~\ref{tab:price-test} 末列。其中最需强调的是 T4：电价与光伏的\emph{池化}相关
\(+0.028\) 几乎为 0——这只是"夜间光伏恒为 0"稀释后的假象，日聚合尺度上的真实结构是 \(-0.343\)；若据此
独立抽样，就会生成"午间光伏大发、同时电价高企"这类不可能情景并\textbf{低估尾部风险}（图~\ref{fig:q4-price}）。

\begin{figure}[htbp]
  \centering
  \includefig[width=0.80\linewidth]{fig08_q4_price.pdf}
  \caption{电价过程的骨架与残差结构（上子图：附件 4 电价的逐槽分位带 10\%--90\% 与 25\%--75\%、
           逐槽中位数与附件 1 分时骨架；下子图：日内残差自相关 \(\mathrm{ACF}(1\text{--}6)\)，
           并标注日级残差自相关与峰谷波动比；范围：2025-01-01 至 12-31 共 365 天 \(\times\) 144 槽；
           单位：电价 元/kWh，自相关为无量纲量；样本：\(n=52\,560\) 个电价观测）}
  \label{fig:q4-price}
\end{figure}

\subsection{两层电价模型与联合情景生成}

本文用\textbf{两层结构}刻画电价：第一层是可预测的\textbf{分时骨架} \(\bar p_t\)（取附件 1 的电价剖面，
全年尺度上等于附件 4 的逐时段均值，最大偏差 \(5.15\times10^{-5}\)~元/kWh），第二层是不可预测的
\textbf{随机残差} \(\varepsilon_{d,t}\)：
\begin{equation}
  p_{d,t}=\bar p_{t}+\varepsilon_{d,t},\qquad
  \mathbb{E}[\varepsilon_{d,t}]=0
  \label{eq:q4-price}
\end{equation}
其中 \(d\) 为日期、\(t\) 为槽位，\(p_{d,t}\) 的单位为\unit{元/kWh}。按 \S7.1 的 T3，残差须
\textbf{按时刻对齐}重采样，故本文直接以\textbf{整日}为单位构造情景。

\textbf{联合情景：二维分层取代表历史日。} 设 \(r_{d,t}\) 为第 \(t\) 槽的净需求预报残差（\(0{:}00\) 预报
口径的净需求与实际净需求之差，单位\unit{kWh}），\(p_{d}\) 为该日平均电价。把 334 天按
\((\sum_t r_{d,t},\ p_d)\) 各分 8 层形成 \(8\times8\) 网格，每格取\textbf{medoid}（层内欧氏距离和最小的
历史日）为代表，权重取层内样本占比。所得 64 个代表日（权重 \(0.0030\)--\(0.0299\)，合计 1）保真两个
维度的边缘分布，并因电价剖面与净需求残差来自\textbf{同一个历史日}而\textbf{天然保留联合结构}。

情景保真度见表~\ref{tab:q4-scenfid}。作为对照，本文另用\textbf{分块 Bootstrap}（6~h 块、按时刻对齐）
生成 64 个合成情景：优势是不依赖历史日的可获取性，但本题样本量下边缘保真度明显更差（均值偏差 3.41\%、
1\% 分位相对差 20.4\%），故\textbf{只作为对照}，主口径采用代表历史日。



\textbf{必须保留 \(\rho(p,L)\) 的理由。} 电价与净负荷的正相关（池化 \(+0.260\)）意味着"贵的时候恰好也在
多买电"；若独立抽样，"高价 \(\times\) 高净需求"的联合尾部占比会被人为压低，风险度量随之失去意义。二维
网格取代表日后 \(\rho(p,\text{净负荷})=0.2597\)，与实际 0.2620 相差 0.0023，尾部结构被保真。

\textbf{波动电价下的 LP 形式与双重不确定性。} 设 \(\Omega\) 为联合情景集（权重 \(w_\omega\)），
\(r_{\omega,t}\) 为情景 \(\omega\) 的净需求残差，\(\widehat{PV}^{(0)}_t\) 为 \(0{:}00\) 发布的光伏预报，
\(p_{\omega,t}\) 为情景电价，则重算问题 2 的两阶段 LP 为
\begin{subequations}\label{eq:q4-lp}
  \begin{equation}
    \min\ \ \sum_{t=0}^{143}\bar p^{\text{sc}}_tE_t
    +\sum_{\omega\in\Omega}w_\omega\sum_{t=0}^{143}5\,p_{\omega,t}\,g_{\omega,t},
    \qquad \bar p^{\text{sc}}_t=\sum_{\omega\in\Omega}w_\omega p_{\omega,t}
    \label{eq:q4-lp-obj}
  \end{equation}
  \begin{equation}
    E_t-c_t+d_t\ \ge\ \bigl(L_t-\widehat{PV}^{(0)}_t\bigr)\Delta t
    \label{eq:q4-lp-plan}
  \end{equation}
  \begin{equation}
    g_{\omega,t}\ \ge\ \bigl(L_t-\widehat{PV}^{(0)}_t\bigr)\Delta t+r_{\omega,t}
    +c_t-d_t-E_t,\qquad g_{\omega,t}\ge 0
    \label{eq:q4-lp-gap}
  \end{equation}
\end{subequations}
约束外加式~\eqref{eq:q3-soc} 的储电量递推、界与终端条件（物理参数同 \S6.4）。两个不确定性分别进入
\textbf{目标系数}（电价 \(p_{\omega,t}\)）与\textbf{约束右端项}（净需求残差 \(r_{\omega,t}\)），即问题 4
相对问题 2 的\textbf{双重不确定性}。决策时的电价信息取情景概率加权均值 \(\bar p^{\text{sc}}_t\)：附件 4
未提供日前电价预报，故 \(0{:}00\) 可用的电价预报取\textbf{分时骨架} \(\bar p_t\)（假设 A10），
\(\bar p^{\text{sc}}_t\) 与骨架的最大偏差 0.0462~元/kWh（占均价 6.02\%）。\textbf{结算一律按附件 4 的当日
实际电价}，使"决策用预报、结算用实现"各归其位（\calref{C-03}）。问题 4-3 在式~\eqref{eq:q4-lp} 上再叠加
\(6{:}00/12{:}00/18{:}00\) 三次滚动修正，调整与违约费用按 \calref{C-09} 的偏差量定价结算。

\subsection{阈值律在波动电价下的形式不变性}

问题 1 的阈值律只依赖①目标对购电量线性、②储电量递推两条结构性质，\textbf{完全不依赖电价的分布假设}，
故电价由确定变随机只改变阈值（水价值 \(\nu_t\)）的\textbf{取值}，不改变"电价低于充电阈值则充、高于放电
阈值则放"的\textbf{函数形式}：\(\alpha_t\le\eta_c\nu_t\) 时充电、\(\alpha_t\ge\nu_t/\eta_d\) 时放电、
\(\eta_c\nu_t<\alpha_t<\nu_t/\eta_d\) 时既不充也不放（死区，两侧之比 \(1/(\eta_c\eta_d)=1.2346\)）。
两点边界：①该结论\textbf{在储电量界与充放功率界不活跃时精确成立}，界活跃时须回到互补松弛式~
\eqref{eq:q1-kkt-c}--\eqref{eq:q1-kkt-d} 并计入界乘子；②"形式不变"不等于"阈值是常数"——波动电价下
\(\nu_t\) 仍须\textbf{逐时段由场景对偶确定}，随情景集与风险参数 \(\lambda\) 变化（\evid{C-Q4-01}）。

\subsection{风险—费用有效前沿}

\textbf{CVaR 的线性化。} 记 \(C_\omega\) 为情景 \(\omega\) 的总购电费用、\(\alpha=0.95\) 为锁定置信水平，
条件风险价值（Rockafellar--Uryasev 形式）为
\begin{equation}
  \mathrm{CVaR}_\alpha(C)=\min_{\eta}\Bigl\{\eta+\frac{1}{(1-\alpha)|\Omega|}
  \sum_{\omega\in\Omega}\max\bigl(0,\,C_\omega-\eta\bigr)\Bigr\},
  \label{eq:q4-cvar}
\end{equation}
引入辅助变量 \(z_\omega\ge C_\omega-\eta\)、\(z_\omega\ge0\) 后，式~\eqref{eq:q4-cvar} 的分段线性结构
完全由线性约束表达，\textbf{加入风险度量后模型仍是 LP}，不引入整数变量。

\begin{caliber}{C-11}{条件风险价值 CVaR（\(\alpha=0.95\)）}
\textbf{定义}\quad 式~\eqref{eq:q4-cvar} 的 \(\mathrm{CVaR}_{0.95}\)，
即最差 5\% 情景下费用的条件期望，按 Rockafellar--Uryasev 线性化求解；
目标函数取凸组合 \(\min\ (1-\lambda)\,\mathbb{E}[C]+\lambda\,\mathrm{CVaR}_{0.95}\)，
\(\lambda\in[0,1]\)（\(\lambda=0\) 为纯期望最优、\(\lambda=1\) 为纯风险最优，
与"\(\min\mathbb{E}[C]+\lambda'\mathrm{CVaR}\)"式在 \(\lambda'=\lambda/(1-\lambda)\) 下等价）；
除假设 A10 外，模型、情景集与物理参数在扫描中保持不变；%
\textbf{数据}\quad 联合情景集（\(8\times8\) 二维分层得到的 64 个代表历史日，权重为层占比）；
评估用 2025-02-01 至 12-31 中\textbf{每 6 天取 1 天}的 56 天子样本，逐日独立求解；%
\textbf{时段}\quad 上述 56 个代表日，单日 144 槽；\textbf{单位}\quad 元/日（日费用的样本均值）、\unit{kWh}/日（电量）；%
\textbf{不含}\quad 不含储能投资与运维成本；也不含"最坏情形"（\(\min\max_\omega C_\omega\)）口径，
故 \(\lambda=1\) 只是\textbf{近似最坏}；%
\textbf{对照}\quad 每一点同时给出四种量：情景内期望费用、情景内 \(\mathrm{CVaR}_{0.95}\)、
\textbf{真实轨迹回测}（在同样的代表日上用附件 4 实际电价与附件 2 实际负载、光伏重算的日费用）与紧急购电量；%
\textbf{复算}\quad \texttt{q4.py --mode cvar --K 64 --step 6}，输出 \texttt{intermediate/q4\_cvar.json} 的 frontier 列表。
\end{caliber}

以 \(\lambda\)（\(0,0.1,\dots,1.0\) 共 11 档）为参数扫描得到\textbf{风险—费用有效前沿}（表~
\ref{tab:frontier-backtest}、图~\ref{fig:frontier}），两端极端解为 \(\lambda=0\) 的纯期望最优与
\(\lambda=1\) 的纯风险最优，后者以 0.99\% 的期望费用换取 3.19\% 的 \(\mathrm{CVaR}\) 下降。



\textbf{推荐点的判据事先锁定}：在"情景内期望费用相对 \(\lambda=0\) 上升不超过 0.5\%"（门限
42\,205.57~元/日）下使 \(\mathrm{CVaR}_{0.95}\) 最小。满足该门限的最大档位是 \(\lambda^{*}=0.3\)：
期望费用 42\,168.43~元/日（\(+0.41\%\)）、\(\mathrm{CVaR}_{0.95}\) 50\,814.83~元/日（\(-2.61\%\)）、
\textbf{真实轨迹回测费用 45\,778.03~元/日（仅 \(+0.18\%\)）}，紧急购电量由 480.19 降到
324.60~\unit{kWh}/日（\(-32.40\%\)）。

\textbf{前沿的拐点。} \(\lambda\) 从 0.7 升到 0.9 时 \(\mathrm{CVaR}_{0.95}\) 只再降 11.53~元/日（0.02\%），
期望费用却继续上升 34.27~元/日（0.08\%）；\(\lambda=1\) 的 \(\mathrm{CVaR}_{0.95}\) 与 \(\lambda=0.9\)
基本持平（50\,509.78 对 50\,510.46~元/日，差 0.68~元/日，属噪声量级）。即\textbf{\(\lambda\approx0.7\)
以后风险已基本降不动}。真实轨迹回测：三个代表点在\textbf{实际发生}的 334 天轨迹口径下日费用为
45\,694.27--45\,952.85~元，情景内的排序在真实轨迹上同样成立但幅度小得多（\(+0.57\%\) vs \(+0.99\%\)）。

\begin{figure}[htbp]
  \centering
  \includefig[width=0.80\linewidth]{fig09_cvar_frontier.pdf}
  \caption{风险—费用有效前沿与真实轨迹回测（左子图：情景期望费用、\(\mathrm{CVaR}_{0.95}\)
           与真实轨迹回测费用随风险厌恶系数 \(\lambda\) 的变化；右子图：期望费用—\(\mathrm{CVaR}_{0.95}\)
           前沿散点，颜色表示 \(\lambda\)，星号标出边际风险下降最快的档位；
           对象：\(\lambda\in\{0,0.1,\dots,1.0\}\) 共 11 档的解；范围：2025-02-01 至 12-31 中每 6 天取 1 天的 56 个代表日；
           单位：费用 元/日；样本：64 个代表情景、11 个 \(\lambda\) 档位，\(\alpha=0.95\)）}
  \label{fig:frontier}
\end{figure}

\subsection{求解结果}

在波动电价下重解问题 2 与问题 3（表~\ref{tab:q4-summary}），两问的费用口径与固定电价下\textbf{完全相同}
（\calref{C-03}、\calref{C-09}），差别只在电价序列取附件 4 实际值、决策电价信息取情景加权均值，故两组
结果可直接比较。



\textbf{电价波动把费用抬高了多少。} 同一 334 天窗口、同一实际负载与光伏轨迹下，波动价下两阶段与四次
滚动结构的费用分别高出 5.81\% 与 5.67\%，无储能基线高出 3.78\%，节省率因此由 13.06\%/15.08\% 降到
11.35\%/13.54\%（分别下降 1.71 与 1.55 个百分点；表~\ref{tab:q4-summary}）。波动的成本上升\textbf{不能由
价格水平解释}：窗口内附件 4 电价均值 0.7575~元/kWh 反而比方时骨架的 0.7662~元/kWh 低 1.13\%，基线费用仍
高 3.78\%——差额来自\textbf{高电价与高净负荷同时出现}（池化相关 \(+0.260\)，见 \S7.2），这正是"必须生成
联合情景"的定量证据。

\textbf{波动电价下的滚动调整依然有效。} 四次预报把波动价下的费用降低 372\,102.66~元（2.47\%），
紧急购电量由 182\,759.86~\unit{kWh} 降到 63\,523.11~\unit{kWh}（降 65.24\%），降幅大于固定价下的
2.33\%——即\textbf{电价越不确定，滚动更新的价值越大}。问题 4-3 的四项费用为：计划相关
13\,979\,692.49~元、违约 224\,406.57~元、超额 270\,998.49~元、紧急购电 247\,055.63~元，合计
14\,722\,153.17~元（1\,472.22~万元），其中违约电量 664\,243.01~\unit{kWh}、超额电量
259\,710.27~\unit{kWh}；违约与超额费用同时高于固定价，说明\textbf{价格不确定使"调整到多少"本身也变得
更难判断}。四次滚动下波动价与固定价的紧急购电量几乎相同（反而少 998.98~\unit{kWh}，\(-1.55\%\)），
而两阶段结构下波动价多买 10\,744.04~\unit{kWh}（\(+6.25\%\)），即\textbf{滚动修正几乎完全吸收了电价波动
带来的额外缺口}。四个指定日期的结果见第八章表~\ref{tab:key4} 与表~\ref{tab:emergency}（逐时段完整结果见
\texttt{result4-2.xlsx}、\texttt{result4-3.xlsx}）。
